\section{Kết quả và phân tích}
\label{sec:results}

\subsection{Kết quả tổng quan}

Bốn cấu hình được đánh giá trên cùng một tập test (fold 1, 873 mẫu).
Bảng~\ref{tab:overall} tổng hợp Accuracy và Macro-F1.

\begin{table}[H]
\centering
\caption{So sánh kết quả tổng quan trên tập test.}
\label{tab:overall}
\begin{tabular}{lcc}
\toprule
\textbf{Cấu hình} & \textbf{Test Accuracy} & \textbf{Test Macro-F1}\\
\midrule
C1 -- CNN baseline            & 0{,}6999 & 0{,}7113\\
C2 -- ResNet18                & 0{,}7434 & 0{,}7625\\
C3 -- ResNet18 + augmentation & 0{,}7274 & 0{,}7398\\
\textbf{C4 -- Ensemble}       & \textbf{0{,}7572} & \textbf{0{,}7759}\\
\bottomrule
\end{tabular}
\end{table}

Bảng trên cho thấy ba điểm chính. Thứ nhất, ResNet18 vượt CNN baseline một cách
rõ ràng -- hơn khoảng $4{,}3$ điểm phần trăm Accuracy và $5{,}1$ điểm Macro-F1.
Thứ hai, augmentation \textit{đơn lẻ} lại không cải thiện mà còn kéo cả hai chỉ
số của ResNet18 xuống đôi chút; lý do nằm ở mức từng lớp và được phân tích kỹ ở
Mục~\ref{ssec:aug}. Thứ ba, ensemble của C2 và C3 cho kết quả tốt nhất ở cả
Accuracy lẫn Macro-F1, và là cấu hình duy nhất vượt mốc $0{,}75$ Accuracy. Như
vậy, dù C3 yếu hơn C2 khi đứng một mình, nó vẫn đóng góp tích cực khi được kết
hợp -- một kết quả thoạt nhìn nghịch lý nhưng hoàn toàn hợp lý, sẽ được giải
thích ở Mục~\ref{ssec:ensemble}.

\subsection{Phân tích cấu hình CNN baseline}

CNN baseline đạt Accuracy $0{,}6999$ và Macro-F1 $0{,}7113$. Đây là một mốc cơ sở
hợp lý: với một kiến trúc bốn khối tích chập đơn giản, mô hình vẫn học được phần
lớn cấu trúc thời gian -- tần số trên log-Mel spectrogram. Điều này gián tiếp xác
nhận pipeline tiền xử lý của đề tài là đúng hướng -- nếu đặc trưng đầu vào kém,
ngay cả mô hình cơ sở cũng khó đạt tới mức này.

\begin{figure}[H]
  \centering
  \imgfallback{wandb_cnn.png}{0.85\textwidth}
  \caption{Đường cong huấn luyện của CNN baseline ghi từ W\&B.}
  \label{fig:wandbcnn}
\end{figure}

F1-score theo từng lớp của CNN (Bảng~\ref{tab:f1cnn}) hé lộ điểm yếu cụ thể. Các
lớp có đặc trưng nổi bật như \texttt{car\_horn} ($0{,}97$) và \texttt{siren}
($0{,}80$) đạt điểm cao. Đáng chú ý nhất là \texttt{engine\_idling}: F1 chỉ
$0{,}37$ nhưng precision tới $0{,}96$ trong khi recall chỉ $0{,}23$. Nói cách
khác, CNN gần như ``ngại'' gán nhãn \texttt{engine\_idling}; khi nó gán thì
thường đúng, nhưng đa số mẫu \texttt{engine\_idling} thực tế lại bị đẩy sang lớp
khác. Đây là dấu hiệu sớm của vấn đề nhóm âm thanh cơ khí mà mọi mô hình về sau
đều phải đối mặt.

\begin{table}[H]
\centering
\caption{F1-score theo lớp -- CNN baseline.}
\label{tab:f1cnn}
\small
\begin{tabular}{lc@{\hspace{2.2em}}lc}
\toprule
\textbf{Lớp} & \textbf{F1} & \textbf{Lớp} & \textbf{F1}\\
\midrule
air\_conditioner  & 0{,}6601 & engine\_idling & 0{,}3697\\
car\_horn         & 0{,}9714 & gun\_shot      & 0{,}7209\\
children\_playing & 0{,}7106 & jackhammer     & 0{,}6306\\
dog\_bark         & 0{,}8360 & siren          & 0{,}8026\\
drilling          & 0{,}6260 & street\_music  & 0{,}7849\\
\bottomrule
\end{tabular}
\end{table}

\subsection{Phân tích cấu hình ResNet18}

ResNet18 đạt Accuracy $0{,}7434$ và Macro-F1 $0{,}7625$, vượt CNN baseline ở cả
hai chỉ số. Kết quả này trả lời trực tiếp câu hỏi \textbf{CH1}: trong điều kiện
thực nghiệm của đề tài, độ sâu lớn hơn cùng kết nối phần dư \textit{có} chuyển
thành hiệu năng tốt hơn. Việc giữ stem nhẹ để bảo toàn chi tiết cục bộ của
spectrogram nhiều khả năng đã góp phần vào kết quả này.

\begin{figure}[H]
  \centering
  \imgfallback{wandb_resnet.png}{0.85\textwidth}
  \caption{Đường cong huấn luyện của ResNet18 ghi từ W\&B.}
  \label{fig:wandbresnet}
\end{figure}

Xét theo lớp (Bảng~\ref{tab:f1resnet}), ResNet18 cải thiện mạnh nhóm cơ khí so
với CNN: \texttt{jackhammer} tăng từ $0{,}63$ lên $0{,}76$, \texttt{engine\_idling}
từ $0{,}37$ lên $0{,}67$. Các lớp dễ như \texttt{gun\_shot} ($0{,}96$) và
\texttt{street\_music} ($0{,}87$) cũng rất tốt. Tuy nhiên, mô hình lại có một
điểm yếu riêng: \texttt{air\_conditioner} rớt xuống F1 $0{,}35$ (recall chỉ
$0{,}31$). ResNet18 dường như thiên về các đặc trưng năng lượng mạnh và bỏ sót
lớp \texttt{air\_conditioner} vốn trầm, đều và ít nổi bật.

\begin{table}[H]
\centering
\caption{F1-score theo lớp -- ResNet18.}
\label{tab:f1resnet}
\small
\begin{tabular}{lc@{\hspace{2.2em}}lc}
\toprule
\textbf{Lớp} & \textbf{F1} & \textbf{Lớp} & \textbf{F1}\\
\midrule
air\_conditioner  & 0{,}3503 & engine\_idling & 0{,}6705\\
car\_horn         & 0{,}9167 & gun\_shot      & 0{,}9577\\
children\_playing & 0{,}7725 & jackhammer     & 0{,}7632\\
dog\_bark         & 0{,}8436 & siren          & 0{,}8449\\
drilling          & 0{,}6400 & street\_music  & 0{,}8660\\
\bottomrule
\end{tabular}
\end{table}

\subsection{Phân tích cấu hình ResNet18 + augmentation}
\label{ssec:aug}

ResNet18 + augmentation đạt Accuracy $0{,}7274$ và Macro-F1 $0{,}7398$ -- cùng
thấp hơn ResNet18 không augmentation một chút. Nếu chỉ nhìn hai con số này, ta dễ
kết luận rằng augmentation không có ích. Nhưng phân tích theo lớp
(Bảng~\ref{tab:f1aug}) cho thấy bức tranh tinh tế hơn nhiều: augmentation không
làm mô hình ``kém đi'' một cách đồng đều, mà \textit{dịch chuyển} năng lực giữa
các lớp.

\begin{table}[H]
\centering
\caption{F1-score theo lớp -- ResNet18 + augmentation.}
\label{tab:f1aug}
\small
\begin{tabular}{lc@{\hspace{2.2em}}lc}
\toprule
\textbf{Lớp} & \textbf{F1} & \textbf{Lớp} & \textbf{F1}\\
\midrule
air\_conditioner  & 0{,}4872 & engine\_idling & 0{,}6749\\
car\_horn         & 0{,}8696 & gun\_shot      & 0{,}8831\\
children\_playing & 0{,}8142 & jackhammer     & 0{,}4725\\
dog\_bark         & 0{,}8585 & siren          & 0{,}8306\\
drilling          & 0{,}6442 & street\_music  & 0{,}8629\\
\bottomrule
\end{tabular}
\end{table}

Một mặt, augmentation cải thiện rõ những lớp ResNet18 vốn yếu hoặc trung bình:
\texttt{air\_conditioner} tăng từ $0{,}35$ lên $0{,}49$, \texttt{children\_playing}
từ $0{,}77$ lên $0{,}81$, \texttt{dog\_bark} từ $0{,}84$ lên $0{,}86$. Việc che
ngẫu nhiên một phần spectrogram buộc mô hình không bám vào một vài chi tiết
mạnh, nhờ đó để mắt hơn tới các lớp trầm.

Mặt khác, đúng nhóm âm thanh cơ khí lại bị tổn thương. \texttt{jackhammer} rớt
mạnh từ $0{,}76$ xuống $0{,}47$, \texttt{gun\_shot} từ $0{,}96$ xuống $0{,}88$.
Nguyên nhân khá trực giác: \texttt{jackhammer} được nhận diện chủ yếu nhờ tính
\textit{lặp đều} theo thời gian; khi time masking che mất một dải thời gian, ta
vô tình xóa luôn chính đặc trưng đó. Vì \texttt{jackhammer} là lớp đông mẫu nhất
trong tập test (120 mẫu), cú rớt này kéo cả Accuracy lẫn Macro-F1 đi xuống. Đây
chính là lý do cường độ augmentation đã phải được giảm bớt (Mục~\ref{sec:method});
ngay cả ở mức nhẹ, sự đánh đổi vẫn còn.

Như vậy, câu trả lời cho \textbf{CH2} là: augmentation không phải một cải tiến
``miễn phí''. Nó đánh đổi độ sắc bén ở nhóm cơ khí lấy độ bao phủ tốt hơn ở nhóm
âm thanh trầm. Đứng một mình, sự đánh đổi này hơi bất lợi -- nhưng nó tạo ra một
mô hình có hồ sơ lỗi \textit{khác hẳn} ResNet18, và đó lại là điều kiện lý tưởng
cho ensemble.

\begin{figure}[H]
  \centering
  \imgfallback{wandb_resnet_aug.png}{0.85\textwidth}
  \caption{Đường cong huấn luyện của ResNet18 + augmentation ghi từ W\&B.}
  \label{fig:wandbaug}
\end{figure}

\subsection{Phân tích cấu hình Ensemble}
\label{ssec:ensemble}

Ensemble lấy trung bình xác suất của ResNet18 và ResNet18 + augmentation, đạt
Accuracy $0{,}7572$ và Macro-F1 $0{,}7759$ -- tốt nhất trong cả bốn cấu hình.
Quan trọng hơn con số, cách ensemble đạt được kết quả này mới là điều đáng phân
tích.

Hai mô hình thành phần mạnh ở hai vùng khác nhau. ResNet18 sắc ở nhóm cơ khí
(\texttt{jackhammer} $0{,}76$, \texttt{gun\_shot} $0{,}96$); ResNet18 +
augmentation bao phủ tốt hơn nhóm trầm (\texttt{air\_conditioner} $0{,}49$,
\texttt{children\_playing} $0{,}81$). Khi trung bình xác suất, mỗi mô hình ``kéo''
phần nó tự tin và làm dịu phần nó yếu. Kết quả là ensemble không đơn thuần là
trung bình cộng của hai mô hình mà thường tốt hơn cả hai ở từng lớp: F1 của
\texttt{jackhammer} hồi về $0{,}67$ (so với $0{,}47$ của C3), \texttt{air\_conditioner}
giữ ở $0{,}42$ (so với $0{,}35$ của C2), còn \texttt{siren} đạt $0{,}90$ -- cao
hơn cả hai mô hình thành phần.

\begin{table}[H]
\centering
\caption{F1-score theo lớp -- Ensemble.}
\label{tab:f1ens}
\small
\begin{tabular}{lc@{\hspace{2.2em}}lc}
\toprule
\textbf{Lớp} & \textbf{F1} & \textbf{Lớp} & \textbf{F1}\\
\midrule
air\_conditioner  & 0{,}4224 & engine\_idling & 0{,}6937\\
car\_horn         & 0{,}9296 & gun\_shot      & 0{,}9444\\
children\_playing & 0{,}8214 & jackhammer     & 0{,}6730\\
dog\_bark         & 0{,}8421 & siren          & 0{,}8989\\
drilling          & 0{,}6569 & street\_music  & 0{,}8763\\
\bottomrule
\end{tabular}
\end{table}

Đây là câu trả lời cho \textbf{CH3}: ensemble \textit{có} khai thác được tính bù
trừ lỗi. Việc một mô hình yếu hơn (C3) vẫn nâng được kết quả chung cho thấy giá
trị của một thành phần ensemble không nằm ở chỗ nó mạnh tuyệt đối, mà ở chỗ nó
\textit{sai khác đi} so với các thành phần còn lại.

\subsection{So sánh trực tiếp bốn cấu hình}

Bảng~\ref{tab:perclass} đặt F1-score theo lớp của cả bốn cấu hình cạnh nhau, giúp
nhìn rõ xu hướng.

\begin{table}[H]
\centering
\caption{F1-score theo lớp của bốn cấu hình. Giá trị in đậm là cao nhất theo hàng.}
\label{tab:perclass}
\small
\begin{tabular}{lcccc}
\toprule
\textbf{Lớp} & \textbf{CNN} & \textbf{ResNet18} & \textbf{ResNet18+aug} & \textbf{Ensemble}\\
\midrule
air\_conditioner  & \textbf{0{,}6601} & 0{,}3503 & 0{,}4872 & 0{,}4224\\
car\_horn         & \textbf{0{,}9714} & 0{,}9167 & 0{,}8696 & 0{,}9296\\
children\_playing & 0{,}7106 & 0{,}7725 & 0{,}8142 & \textbf{0{,}8214}\\
dog\_bark         & 0{,}8360 & 0{,}8436 & \textbf{0{,}8585} & 0{,}8421\\
drilling          & 0{,}6260 & 0{,}6400 & 0{,}6442 & \textbf{0{,}6569}\\
engine\_idling    & 0{,}3697 & 0{,}6705 & 0{,}6749 & \textbf{0{,}6937}\\
gun\_shot         & 0{,}7209 & \textbf{0{,}9577} & 0{,}8831 & 0{,}9444\\
jackhammer        & 0{,}6306 & \textbf{0{,}7632} & 0{,}4725 & 0{,}6730\\
siren             & 0{,}8026 & 0{,}8449 & 0{,}8306 & \textbf{0{,}8989}\\
street\_music     & 0{,}7849 & 0{,}8660 & 0{,}8629 & \textbf{0{,}8763}\\
\midrule
\textbf{Macro-F1} & 0{,}7113 & 0{,}7625 & 0{,}7398 & \textbf{0{,}7759}\\
\bottomrule
\end{tabular}
\end{table}

Bảng này làm rõ vài điều. Ensemble chiếm vị trí cao nhất ở sáu trên mười lớp và
ở Macro-F1 tổng -- nó không phải mô hình giỏi nhất ở mọi lớp, nhưng ổn định nhất.
\texttt{air\_conditioner} là một ngoại lệ thú vị: CNN baseline lại đạt F1 cao
nhất ($0{,}66$) cho lớp này, trong khi mọi mô hình ResNet đều chật vật. Điều đó
gợi ý rằng \texttt{air\_conditioner} không phải lớp ``khó tuyệt đối''; ResNet
đơn giản là đã học một biểu diễn ưu tiên các lớp năng lượng mạnh và đánh đổi bằng
lớp trầm này. Cuối cùng, \texttt{jackhammer} cho thấy rõ tác động hai chiều của
augmentation: cao nhất ở ResNet18 ($0{,}76$), thấp hẳn ở ResNet18+aug ($0{,}47$),
rồi được ensemble kéo lại mức an toàn ($0{,}67$).

\subsection{Phân tích confusion matrix}

Confusion matrix của bốn cấu hình (Hình~\ref{fig:cm_cnn}--\ref{fig:cm_ens}) xác
nhận chẩn đoán đã nêu: gần như toàn bộ sai số tập trung trong nhóm âm thanh cơ
khí và âm nền liên tục. Các cặp nhầm lẫn lặp đi lặp lại gồm:

\begin{itemize}[leftmargin=1.4em,itemsep=1pt]
  \item \texttt{air\_conditioner} $\leftrightarrow$ \texttt{engine\_idling} --
        hai âm nền liên tục, phổ năng lượng rất gần nhau;
  \item \texttt{drilling} $\leftrightarrow$ \texttt{jackhammer} -- hai âm cơ khí
        có dải tần mạnh tương tự;
  \item các mẫu âm nền yếu bị đẩy sang lớp cơ khí hoặc \texttt{children\_playing}.
\end{itemize}

Confusion matrix cũng cho thấy ensemble không tạo ra một cơ chế mới để tách nhóm
cơ khí, mà chỉ làm \textit{mờ bớt} các ô nhầm lẫn nhờ trung bình hóa hai quan
điểm. Nói cách khác, ranh giới khó nhất của bài toán -- nhóm machinery -- vẫn còn
nguyên đó; ensemble giảm thiểu hậu quả chứ chưa giải quyết tận gốc.

\begin{figure}[H]
  \centering
  \imgfallback{cnn_baseline_confusion_matrix.png}{0.66\textwidth}
  \caption{Confusion matrix -- CNN baseline.}
  \label{fig:cm_cnn}
\end{figure}

\begin{figure}[H]
  \centering
  \imgfallback{resnet18_regular_v3_confusion_matrix_tta.png}{0.66\textwidth}
  \caption{Confusion matrix -- ResNet18.}
  \label{fig:cm_resnet}
\end{figure}

\begin{figure}[H]
  \centering
  \imgfallback{resnet18_aug_confusion_matrix_tta.png}{0.66\textwidth}
  \caption{Confusion matrix -- ResNet18 + augmentation.}
  \label{fig:cm_aug}
\end{figure}

\begin{figure}[H]
  \centering
  \imgfallback{ensemble_confusion_matrix_tta.png}{0.66\textwidth}
  \caption{Confusion matrix -- Ensemble.}
  \label{fig:cm_ens}
\end{figure}

\subsection{Tổng hợp trả lời câu hỏi nghiên cứu}

\textbf{CH1 -- ResNet18 so với CNN baseline.} ResNet18 vượt CNN ở cả Accuracy
($0{,}74$ so với $0{,}70$) lẫn Macro-F1 ($0{,}76$ so với $0{,}71$). Độ sâu cùng
kết nối phần dư, kết hợp với stem được giữ nhẹ, đã thực sự mang lại lợi ích trên
bài toán này.

\textbf{CH2 -- Tác động của augmentation.} Augmentation không cải thiện hiệu năng
tổng thể khi dùng đơn lẻ. Tác động thật sự của nó là \textit{tái phân bố} năng
lực: cải thiện các lớp trầm như \texttt{air\_conditioner}, nhưng làm giảm độ sắc
ở nhóm cơ khí, đặc biệt \texttt{jackhammer}. Đây là một sự đánh đổi chứ không
phải một cải tiến thuần.

\textbf{CH3 -- Giá trị của ensemble.} Ensemble cho kết quả tốt nhất, và là cấu
hình duy nhất vượt $0{,}75$ Accuracy. Nó thành công không phải vì ghép hai mô
hình mạnh, mà vì ghép hai mô hình có \textit{lỗi khác nhau} -- minh chứng cho
nguyên tắc cốt lõi của ensemble learning.

\subsection{Hạn chế của đề tài}

Dù pipeline đã hoàn chỉnh và kết quả nhất quán, đề tài vẫn còn một số hạn chế cần
nêu thẳng thắn:

\begin{itemize}[leftmargin=1.4em,itemsep=1pt]
  \item Mọi kết quả dựa trên \textit{một} cách chia fold (test = fold 1). Giao
        thức chuẩn của UrbanSound8K là kiểm định chéo 10-fold; con số trung bình
        trên 10 fold sẽ đáng tin hơn và cho biết cả độ dao động.
  \item Nhóm âm thanh cơ khí vẫn là nút thắt chưa được giải quyết tận gốc; đề
        tài mới giảm nhẹ hậu quả qua ensemble.
  \item Chưa khảo sát có hệ thống không gian siêu tham số, cũng như chưa thử các
        biểu diễn đặc trưng khác (ví dụ thêm delta, delta-delta) có thể giúp tách
        nhóm cơ khí.
  \item Chưa dùng mô hình tiền huấn luyện quy mô lớn; đây là một hướng rõ ràng để
        nâng trần hiệu năng.
\end{itemize}

Trong phạm vi một bài tập lớn cá nhân với tài nguyên hạn chế, các kết quả hiện có
vẫn đủ để rút ra những kết luận có ý nghĩa, trình bày ở mục tiếp theo.
